\scanpage{230}

\inserted{able} to introduce \deleted{(at least for \unsure{dim $d$|dim $2$}} on the group of divisor-classes of $A$ a structure of an analytic group (which in fact is algebraic\ldots), which from the algebraic point of view should be possible rather for the group of divisor classes of the completion $\widehat{A}$; of course Mumford uses directly the same kind of argument I used.

I do not know if in the last statement, the hypothesis that $Y'$ is regular (``$y$ isolated singularity'') is essential; we could dispense with it (\deleted{if you} and replace it simply by ``$A$ reduced'' if you can prove by your theory of resolution the following: if $f:X\rightarrow Y$ is proper ``birational'', $X$ regular, then $R^1f_*(\underline{O}_X)=0$. I understand you proved this if $Y$ also is regular (which is easily checked by your theory), but I wonder if this \uline{is really} needed.\margin{of course it is} I would not be surprised either if in this statement, $Y'$ can be replaced by any open subset of $Y$ (replacing of course $\widehat{Y}'$ by the inverse image of the latter). Moreover, I would expect the analogous statements to hold for $\pi_1$, more generally for all ``topological'' invariants as Weil homology, homotopy groups etc, that can be defined for schemes. \deleted{This}\inserted{This} should be related to the fact that all these invariants vanish for the geometric fibers of the morphism $\operatorname{Spec}(\widehat{A})\quad\operatorname{Spec}(A)$. This is easy to check at least for $\pi_0$ (and is true in fact for any henselian ring which is a ``good'' ring); however I do not know if this is true also for $\pi_1$.

Besides, I would not be surprised if most of the previous results (namely parts (ii) and (iii) of the second theorem, and the two corollaries, as well as the \deleted{previous} conjectures of the previous section) did hold true for any ``good'' ring which is henselian, \inserted{or} at least for the

\mspage{229}
